%#####################################################
\section{Introduction}
\label{sec:intro}
%#####################################################
Weak gravitational lensing (WL) causes subtle distortion of observed galaxy shapes due to the gravitational influence of intervening large-scale structures along the line of sight. By statistically analyzing correlations in these shear distortions across large samples of galaxies, weak lensing surveys provide a powerful probe of the large-scale structure of the universe \citep{Bartelmann_2001, Kilbinger_2015}. These analyses allow us to constrain key cosmological parameters, particularly $\Omega_m$ (matter density) and $\sigma_8$ (amplitude of matter fluctuations), test the standard $\Lambda$CDM model and explore potential deviations from it. WL also causes distortions of cosmic microwave background, with the latest constraints presented in \cite{Qu2025_ACTSPTPLANCK_CMBL_S8}. WL cross-correlation with galaxies can also be combined with galaxy auto-correlation for additional 
constraints on cosmological parameters.
This technique however suffers from potential systematics in the modeling of the cross-correlation coefficient between galaxies and dark matter \cite{Guzik01}, as well as from 
observational systematics 
 in the galaxy auto-correlation. In 
this paper we focus on WL analyses without 
galaxy clustering. 

Several Stage-III weak lensing surveys, including the Hyper Suprime-Cam (HSC; \cite{Dalal2023CosmoShearHSC, Li2023_HSCY3_CosmicShear}), the Kilo-Degree Survey (KiDS; \cite{Wright2025_KiDSLegacy}), and the Dark Energy Survey (DES; \cite{DESY3Amon,DESY3Secco}), have made significant progress in constraining the cosmological model. However, a persistent feature of these studies has been the so-called $S_8$ tension: weak lensing surveys have consistently reported values for the lensing amplitude $S_8\equiv\sigma_8\sqrt{\Omega_m/0.3}$ that are systematically lower than those inferred from cosmic microwave background (CMB) measurements \cite{Planck2018_CMB} under the $\Lambda$CDM framework. This discrepancy has led to substantial interest, as it may point to either unknown systematics in the analyses or new physics beyond the standard model. 

Disentangling these possibilities requires improved control of systematics, particularly in the context of next-generation weak lensing surveys, such as the Rubin Observatory Legacy Survey of Space and Time (LSST) \citep{LSST2019survey}, the Nancy Grace Roman Space Telescope \citep{roman2019wfirst}, and Euclid \citep{EuclidOverview}. These surveys will achieve unprecedented precision, making the mitigation of systematics a crucial priority. Among various systematics, photometric redshift calibration is particularly critical, as errors in redshift distributions propagate directly into biases in the constraints of $S_8$.

%Recent advancements in weak lensing surveys have demonstrated the importance of continually refining analysis pipelines to reduce systematic uncertainties and improve the robustness of cosmological constraints. For example, 
Most recently, the KiDS Legacy survey introduced several enhancements to its analysis, including improved photometric redshift calibration, more precise covariance estimation, and updated intrinsic alignment modeling \citep{Wright2025_KiDSLegacy}. Among these, improvements in photometric redshift calibration played a particularly significant role, leading to tighter constraints on cosmological parameters and a notable shift in their $S_8$ results. The Hyper Suprime-Cam (HSC) Year 3 (Y3) analysis, on the other hand, adopted a wide prior on photometric redshift biases due to limited photo-z calibration in higher redshift bins \citep{Li2023_HSCY3_CosmicShear, Dalal2023CosmoShearHSC}. This approach allowed HSC to account for uncertainties in the redshift distributions and enabled robust inference, but it also resulted in weaker constraints on $S_8$. Both KiDS and HSC highlight the critical role of photometric redshift calibration in reducing systematic uncertainties and improving the precision of weak lensing cosmology.

In this work, we reanalyze the HSC Y3 shear-shear correlation function using the updated clustering redshift calibration presented in \cite{ChoppinDeJanvry2025a}, which leverages spectroscopic data from the Dark Energy Spectroscopic Instrument (DESI, \cite{DESI2016b.Instr, DESI2022.KP1.Instr}). Our goal is to quantify the impact of this improved redshift calibration on the cosmological constraints derived from HSC Y3 and to explore its implications for the $S_8$ tension. Our analysis demonstrates that the improved photometric redshift calibration significantly reduces systematic uncertainties, yielding tighter and more robust constraints. Importantly, we find that the updated HSC Y3 results are consistent with Planck CMB measurements, providing a potential resolution to the $S_8$ tension within the context of HSC. We also compare the updated HSC Y3 constraints with those from KiDS Legacy and DES Y3, examining the consistency among surveys. Additionally, we demonstrate the power of combining galaxy lensing and CMB lensing constraints, achieving sub-percent precision comparable to the Planck $\times$ ACT $\times$ DESI results.

This paper is organized as follows. In Section \ref{sec:data}, we describe the dataset used in this paper, including HSC Y3 shape catalog and the galaxy redshift distribution $n(z)$ calibrated with the DESI data. In Section \ref{sec:method}, we present the methodology followed in this analysis to recover cosmological constraints. In Section \ref{sec:results}, we present the cosmological constraints and compare to other experiments and surveys. Finally we conclude in Section \ref{sec:conclusion}.

%%\todo{What to keep/what to leave from this other intro}

%%The distortion of light from distant galaxies caused by the gravitational influence of intervening large-scale structures that trace total matter in the universe is called Weak Gravitational Lensing (WL). This phenomenon creates a subtle cosmic shear pattern in the sky, altering the observed shapes and orientations of background galaxies. The distortion of galaxy shapes, quantified through summary statistics, holds valuable information about underlying cosmological parameters \citep{Bartelmann_2001, Kilbinger_2015}. 
%%With the current and upcoming generations of large-scale optical imaging surveys, such as the Hyper Suprime-Cam on Subaru (HSC, \cite{HSC2022datarelease3}), the Dark Energy Survey (DES, \cite{DES2018datarelease1}), the Kilo-Degree Survey (KiDS, \cite{KiDS2025datarelease5}), the Nancy Grace Roman Space Telescope \cite{roman2019wfirst}, Euclid \cite{EuclidI2024overview} and the Vera C. Rubin Observatory's Legacy Survey of Space and Time (LSST, \cite{LSST2019survey}) probing hundreds of millions to billions of galaxies with accurate shape information, the strength of cosmological parameter inference from weak lensing analysis is becoming increasingly important.
%%These surveys provide or will provide observational data that can be used to constrain fundamental cosmological parameters, particularly $\Omega_m$ (matter density) and $\sigma_8$ (amplitude of matter fluctuations) that WL signals are most sensitive to in the standard cosmological model \citep{Joudaki_2016, Kohlinger_2017, Hikage_2019, Hamana_2020}.

%%To extract information from WL data, various summary statistics are employed. At the two-point level, WL data are analyzed using shear-shear correlation functions or the power spectrum of the shear or convergence. 
%%While these traditional summary statistics leave a wealth of information untapped in the WL signal due to the highly non-Gaussian features at small scales, they are the simplest to extract from the data and compare to the theoretical models. 
%%Still, some of the most precise single measurements of the lensing amplitude $S_8\equiv\sigma_8\sqrt{\Omega_m/0.3}$ currently come from cosmic shear analyses \cite{CosmoShearLi2023HSCY3, Dalal2023CosmoShearHSC, Amon2022CosmoShearDES, Secco2022CosmoShearDES, Asgari2021KidsCosmoShear, wright2025kidslegacyCosmoShear}.

%%However, in order to perform this measurement, optical surveys must rely on their normalized sample redshift distributions projected along the line of sight $n(z)$ when measuring two-point statistics of gravitational shear fields. 
%%Since optical surveys resort to broadband photometry in few optical and near-infrared filters, accurate redshift information cannot be derived from emission lines, which get smoothed out by the broadband nature of the filters. Therefore, imaging surveys probe the $n(z)$ of their dataset through photometric redshift (photo-$z$) algorithms, though these algorithms are usually orders of magnitude less precise than spectroscopically measured redshifts, and easily subject to potential bias or outliers. Nonetheless, the photometric redshift still allow to select galaxies in different tomographic bins defined by photo-$z$ bounds : the $n(z)$ is therefore estimated for each bin. In the case of the fiducial HSC Year 3 analysis \cite{HSCClusteringRau2023, CosmoShearLi2023HSCY3, Dalal2023CosmoShearHSC}, four tomographic bins are used, spanning linear bins from $z_p\sim0.3$ to $z_p\sim1.5$, where $z_p$ is the best photometric redshift estimate of the \texttt{DNNz} method \cite{DNNZHSCNishizawa2020}. 

%%Because of the inaccuracies of photo-$z$ algorithms, the underlying true distribution needs to be calibrated with external information than just the photometric redshifts. One of the common approaches involve using accurate spectroscopic redshifts with color information in order to calibrate the color-redshift relation of the galaxy population. Such a technique is often executed with machine learning methods such as Self-Organizing Maps (SOM). In the companion study \cite{ChoppinDeJanvry2025a}, we leverage the projected galaxy cross-correlation clustering signal with narrow spectroscopic bins to infer the $n(z)$ of the tomographic bins of HSC. This method, also known as clustering redshifts, has been explored in some of the following works 
% maybe reduce number of citations here
%%\cite{Newman2008ClusterZ,ménard2014clusteringbasedredshiftestimationmethod, Gatti2018DESY1ClusteringRedshifts, Gatti2018DESY1systematics,Gatti2021ClusteringDESWLsrcDistrib, Cawthon2022DESY3Boss,Davis2017DESyear1ClusterZWLsrcdistrib,EuclidClusterRedshifts2025,KidsCrossCorrvdB2020,wright2025kidslegacyClusterZ,HSCClusteringRau2023,Morrison2017TheWiZZClusterZ, Schmidt2013ClusteringRedshifts}.

%The limiting factor of this approach is often the source depth and color-space coverage of spectroscopic surveys, as spectroscopic redshift targets are more prone to failures and need longer integration times on faint targets with few easily discernible features for spectroscopy. 

%Various surveys, such as the \href{https://www.darkenergysurvey.org/}{Dark Energy Survey} (DES), KIDS \cite{}, \href{https://hsc.mtk.nao.ac.jp/ssp/}{Hyper Suprime-Cam Survey} (HSC), \href{https://sci.esa.int/web/Euclid}{Euclid}, the \href{https://lsst.org}{Vera Rubin Observatory} (Rubin), and the \href{https://roman.gsfc.nasa.gov/}{Nancy Grace Roman Space Telescope} (Roman), aim to map this cosmic shear across large areas of the sky. 
%With the current and upcoming generations of large-scale optical imaging surveys, such as the Hyper Suprime-Cam Subaru Strategic Program (HSC-SSP, \cite{HSC2022datarelease3}), the Dark Energy Survey (DES, \cite{DES2018datarelease1}), the Kilo-Degree Survey (KiDS, \cite{KiDS2025datarelease5}), the Nancy Grace Roman Space Telescope \cite{roman2019wfirst}, Euclid \cite{EuclidI2024overview} and the Vera C. Rubin Observatory's Legacy Survey of Space and Time (LSST, \cite{LSST2019survey}) probing hundreds of millions of galaxies with accurate shape information, the strength of cosmological parameter inference from weak lensing analysis is becoming increasingly important. Some of the most precise single measurements of the lensing amplitude $S_8\equiv\sigma_8\sqrt{\Omega_m/0.3}$ currently come from cosmic shear analyses \cite{CosmoShearLi2023HSCY3, Dalal2023CosmoShearHSC, Amon2022CosmoShearDES, Secco2022CosmoShearDES, Asgari2021KidsCosmoShear, wright2025kidslegacyCosmoShear}, where one measures the correlation in the distortion of galaxy images due to weak gravitational lensing. 

%Since optical surveys rely on broadband photometry covering optical and near-infrared wavelengths, narrow spectral galaxy features such as emission lines are difficult to identify due to the limited coverage per observed object and the broadband nature of the optical filters. Therefore, imaging surveys employ photometric redshift (photo-$z$) algorithms to infer individual redshifts. Typically, photometric redshifts are used to split the galaxies into tomographic redshift bins, and the true redshift distribution along the line of sight needs to be estimated for each tomographic bin. 

%One can reconstruct the overall redshift distribution $n(z)$ from observed data for each tomographic bin. In the case of the fiducial HSC Year 3 analysis \cite{HSCClusteringRau2023, CosmoShearLi2023HSCY3, Dalal2023CosmoShearHSC}, four tomographic bins are used, spanning linear bins from $z_p\sim0.3$ to $z_p\sim1.5$, where $z_p$ is a photometric redshift from a given photo-$z$ method. Here forth, Bin 1 refers to $0.3\leq z_p<0.6$, Bin 2 to $0.6\leq z_p<0.9$, Bin 3 to $0.9\leq z_p<1.2$ and Bin 4 to $1.2\leq z_p<1.5$ respectively. 

%Many different algorithms have been developed for photometric redshifts, ranging from inferring redshifts based on prior knowledge of galaxy spectral energy distributions (SEDs) (\cite{LePhare1, LePhare2, Brammer2008EAZY, MizukiHSCTanaka2015}) to mapping redshifts to photometry information with machine learning (\cite{CarrascoKind2014PhotozMLz, Hsieh2014DEMp, Collister2004ANNz}). Often, since the performance of the algorithms does not allow for perfect point estimates of the redshift, photo-$z$ methods will instead provide a probability density function (PDF) for each source, which allows for more information on how likely the point estimate is for a given photo-$z$ algorithm. If photometric redshifts were perfectly calibrated, the convolution of these PDFs would accurately describe the PDF $n(z)$ for the tomographic bins. This is not possible in practice due to the performance limitations of the photo-$z$ algorithms. 

%Thus, challenges arise when debiasing and calibrating the redshift distributions obtained by the chosen photo-$z$ algorithm(s), and as a result, weak lensing cosmological analyses can suffer from this systematic. Differences in tomographic redshift calibration can lead to significant cosmological parameter shifts, as demonstrated by Figure 1 of \cite{HSCClusteringRau2023} or by the differences in calibrations on the KiDS-Legacy dataset \cite{wright2025kidslegacyClusterZ, wright2025kidslegacyCosmoShear}, compared to the analysis performed on KiDS-1000 \cite{Asgari2021KidsCosmoShear}. 

%An alternative approach was taken by the HSC team to parametrize photometric redshift uncertainty using a free shift parameter of the redshift first moment for the higher redshift tomographic bins lacking good redshift calibration. While this can encompass systematics, it weakens the obtained cosmological information on $S_8$ and $\Omega_m$ \cite{CosmoShearLi2023HSCY3, Dalal2023CosmoShearHSC}. 

%To mitigate the weaknesses of photo-$z$ algorithms and to debias the photo-$z$ distributions, a number of methods have been developed to calibrate the sample redshift distribution of each bin. 
%A popular approach is to rely on the optical photometry information given by the survey and known spectroscopic redshifts from external datasets to infer redshift distributions from the galaxy population color-space, often using data reduction algorithms and machine learning methods. This technique is commonly executed with Self-Organizing Maps (SOM) \cite{WrightSOMz, SOMPZDES2024, Hildebrandt2021ClusteringRedshiftsSOM}, or through direct weighted calibration \cite{LimaDIRMethod2008}. 
%The limiting factor of this approach is often the source depth and color-space coverage of spectroscopic surveys, as spectroscopic redshift targets are more prone to failures and need longer integration times on faint targets with few easily discernible features for spectroscopy. 

%Another method is to use the clustering of galaxies: by leveraging angular cross-correlations between a spectroscopic sample in small redshift slices and the galaxies in each of the tomographic bins of the weak lensing survey dataset, the technique known as clustering redshifts allows to debias \cite{Newman2008ClusterZ,ménard2014clusteringbasedredshiftestimationmethod, Gatti2018DESY1ClusteringRedshifts, Gatti2018DESY1systematics,Gatti2021ClusteringDESWLsrcDistrib, Cawthon2022DESY3Boss,Davis2017DESyear1ClusterZWLsrcdistrib,EuclidClusterRedshifts2025,KidsCrossCorrvdB2020,wright2025kidslegacyClusterZ,HSCClusteringRau2023,Morrison2017TheWiZZClusterZ, Schmidt2013ClusteringRedshifts} the $n(z)$ distributions. 
%This is the method employed in this work. 
%However, clustering redshifts is also subject to systematics, since the method needs information from small angular scales (typically $0.1-5\hMpc$), and galaxy bias modeling at these scales remains difficult. 
%In this work, common possible systematics that can affect and bias the clustering redshifts calibration of the photometric sample are investigated, such as the impacts of galaxy bias and magnification. 
%This methodology is applied to the S19A Shape Catalog \cite{HSCShapeCataloguePDR3Li2022} from the Hyper Suprime-Cam Subaru Strategic Survey Public Data Release 3 \cite{HSC2022datarelease3} (S19A HSC-SSP PDR3) as our photometric catalog to calibrate. 
%Hereafter, we refer to this dataset simply as HSC. 
%The Dark Energy Spectroscopic Instrument Data Release 1 \cite{DESI.DR1.I.Presentation} and Data Release 2 \cite{DESI.DR2.Presentation.InPrep} (DESI DR1, DR2) will serve as our spectroscopic calibration samples.   

%The Dark Energy Spectroscopic Instrument (DESI) is a Stage IV spectroscopic survey \cite{DESI2016b.Instr, DESI2022.KP1.Instr, Corrector.Miller.2023} carrying out a eight year program that started collecting data in 2021. DESI uses a focal plane hosting 5000 robotic positioners containing optical fibers \cite{FiberSystem.Poppett.2024} to obtain spectra, and therefore accurate spectroscopic redshifts \cite{Spectro.Pipeline.Guy.2023}, from galaxies. DESI plans to probe $17,000$ deg$^2$ of the sky in its main phase of operations \cite{SurveyOps.Schlafly.2023}, covering spectroscopic redshift ranges up to $z\sim4.2$. The DESI experiment studies six different types of tracers: the Milky Way Survey, observing stars in our local galaxy, the Bright Galaxy Sample (BGS, $0.01\lesssim z\lesssim0.6$), the Luminous Red Galaxies (LRG, $0.3\lesssim z\lesssim1.2$), the Emission Line Galaxies (ELG, $0.7\lesssim z\lesssim1.6$), the Quasi-Stellar Objects (QSO, $0.7\lesssim z\lesssim2.1$) and the Lyman-$\alpha$ forest (Ly$\alpha$, $z\gtrsim 2.1$). With its accurate redshift catalog, DESI has already allowed for significant breakthroughs on cosmological measurements \cite{DESI.DR2.II.BAO, DESI2024.VII.KP7B}.

%The Hyper Suprime-Cam Subaru Strategic Program is a Stage 3 optical imaging survey using the Hyper Suprime-Cam wide-field  camera installed on the Subaru 8.2m telescope on the summit of Maunakea, Hawaï, and covers three depth layers: the Wide field, the Deep field, and Ultra-Deep field. The S19A (from the the internal September 2019 data release) Shape Catalog \cite{HSCShapeCataloguePDR3Li2022} compiles galaxy shapes from $i_{\mathrm{band}}$ imaging data obtained from 2014 to 2019 in the Wide field of the HSC-SSP. It is crucial to obtain well calibrated redshift distributions of this dataset, as constraints from cosmic shear heavily depend on these.

%In the tomographic bin analysis performed by the HSC collaboration \cite{HSCClusteringRau2023}, there was no calibration from angular clustering of galaxies above $z\gtrsim1.2$ due to the lack of available spectroscopic datasets, therefore not providing calibration information from clustering redshifts for most of Bin 3 and none for Bin 4. 
%In turn, this makes the mean value of a given redshift bin uncertain, meaning that major shifts are possible on Bins 3 and 4.
%Indeed, analysis by \cite{Zhang2022PhotometricRedshiftShifts} showed a single shift parameter was sufficient to capture $n(z)$ uncertainty per bin in the HSC distributions.
%The cosmic shear analyses on the two-point correlation function \cite{CosmoShearLi2023HSCY3} and on the angular power spectrum \cite{Dalal2023CosmoShearHSC} of HSC used a conservative uniform prior $\mathcal{U}([-1,1])$ for the tomographic bins's redshift shift parameters.
%Both analyses report hints of photo-$z$ miscalibration, with possible redshift expectation values being higher than initially obtained from the redshift bins only calibrated with photometry. Reported marginalized mode of shifts with asymmetric $34\%$ confidence interval and maximum a posteriori (MAP) values of shifts obtained by both HSC cosmic shear analyses have returned\footnote{Note the following convention: a negative shift implies the found expectation value (from cosmic shear analysis) for a bin is higher than the available measurements.}: %$\dz=\langle z\rangle_{\mathrm{meas}} - \langle z\rangle_{\mathrm{true}}$
%$\dz^{(3)}=-0.075^{+0.056}_{
%-0.059} (-0.046)$ and $\dz^{(4)}=-0.157^{+0.094}_{-0.111} (-0.144)$ 
%for the angular power spectrum \cite{Dalal2023CosmoShearHSC} and 
%$\dz^{(3)}=-0.115^{+0.052}
%_{-0.058} (-0.120)$ and $
%\dz^{(4)}=-0.192^{+0.088}
%_{-0.088} (-0.190)$ 
%for the angular two point correlation functions \cite{CosmoShearLi2023HSCY3} in Bin 3 and Bin 4. 
%More recent works \cite{Rana2025HSCY3CosmicShearRatios} used Cosmic Shear Ratios in order to provide calibration for the photometric redshift bins. 
%While the results found by that calibration are consistent with those obtained by the cosmic shear analyses, the large error bars loosely constrain the redshift expectation: the obtained results were $\Delta z_2 = +0.002^{+0.021}_{-0.022},\;\Delta z_3 = -0.002^{+0.085}_{-0.217},\;\Delta z_4 = -0.292^{+0.229}_{-0.324}$. In this work, the focus will be on constraining these two shift parameters, comparing results to the HSC analyses as well as previous calibrations, and producing $n(z)$ distributions for each of the bins. 

%The outline of this paper is as follows. Section \ref{sec:data} presents the datasets employed in this study: HSC’s Shape catalog and DESI DR1, DR2, as well as the quality cuts applied to both catalogs. Section \ref{sec:modeling} describes the modeling of the $n(z)$ distributions with two-point correlation functions and discusses systematic effects, including galaxy bias from both samples and magnification effects. Section \ref{sec:2pcf} describes our measurements of the angular correlation vectors. Section \ref{sec:parametrization} presents a spline based approach to model the $n(z)$ distributions and discusses other methods to regularize the clustering redshifts measurements. Finally, results are showcased in section \ref{sec:results}.